In all the quicksort implementations in this article, we use the
``median-of-three'' pivot selection technique to prevent pivot skewness
\code{part}\cite{sedgewick1977,sedgewick:1978}.
Correspondingly, we ensure that the array length is at least $3$ to call function \code{part}.

\begin{table}[b]
    \centering
    \caption{\label{tab:quicksort-experiment}
    Runtime measurement of Hoare's quicksort algorithms,
    comparing the traditional implementation and the single-loop implementation.
    Experiments are run against three data sets---the first is sorted ascending,
    the second descending, and the third is randomly ordered.
    Courtesy: The experiment is conducted based on project ``Benchmarking Sorting
    Algorithms''\cite{benchmarking-sorting-algo} with some modification.}
    \vskip 10pt
    \begin{tabular}{r|r|r|r}
        \hline
        \multicolumn{4}{l} {\mbox{Integer arrays in ascending order}}\\
        array size & Traditional & Simplified & Percent difference\\
        \hline
        $10$ & $7$ & $8$ & $14\%$\\
        $100$ & $18$ & $20$ & $11\%$\\
        $1,000$ & $111$ & $126$ & $14\%$\\
        $10,000$ & $1,071$ & $1,239$ & $16\%$\\
        $50,000$ & $4,893$ & $5,867$ & $20\%$\\
        $100,000$ & $9,479$ & $11,131$ & $17\%$\\
        $1,000,000$ & $115,248$ & $137,658$ & $19\%$\\
        \hline
        \multicolumn{4}{l} {\mbox{Integer arrays in descending order}}\\
        array size & Traditional & Simplified & Percent difference\\
        \hline
        $10$ & $7$ & $7$ & $0\%$\\
        $100$ & $18$ & $19$ & $6\%$\\
        $1,000$ & $106$ & $142$ & $34\%$\\
        $10,000$ & $1,020$ & $1,167$ & $14\%$\\
        $50,000$ & $4,757$ & $5,411$ & $14\%$\\
        $100,000$ & $9,378$ & $11,132$ & $19\%$\\
        $1,000,000$ & $118,696$ & $138,311$ & $17\%$\\
        \hline
        \multicolumn{4}{l} {\mbox{Integer arrays randomly ordered}}\\
        array size & Traditional & Simplified & Percent difference\\
        \hline
        $10$ & $6$ & $6$ & $0\%$\\
        $100$ & $25$ & $25$ & $0\%$\\
        $1,000$ & $264$ & $249$ & $-6\%$\\
        $10,000$ & $2,949$ & $3,065$ & $4\%$\\
        $50,000$ & $16,565$ & $17,943$ & $8\%$\\
        $100,000$ & $34,005$ & $36,400$ & $7\%$\\
        $1,000,000$ & $380,477$ & $406,659$ & $7\%$\\
        \hline
    \end{tabular}
\end{table}

~\autoref{tab:quicksort-experiment} shows the runtime analysis and comparisons.
The simplified implementation indeed comes with an overhead and is thus slower than
its traditional counterpart of the Hoare's quicksort program.
For sorted data sets (either ascending or descending), there is a $17\%$ slowdown.
But for randomly ordered data sets, the slow down is consistantly around $7\%$.
The experiment data seems to indicate that the simplified
implementation of Hoare's quicksort algorithm shares same runtime complexity as its
traditional counterpart.
Their performance may differ by a multiplier.
